\section{第一轮独立评测、消融与失败分析}\label{sec:results}

\subsection{常规测试的总体误差与对照}

本次基准累计完成5040次定位预测，壁钟时间为0.8分钟。主测试360例采用三个独立数据种子，各常规场景样本数相同。所有方法使用同一案例；所选配置在评测前已冻结。

\begin{table}[H]\centering\small
\caption{独立主测试总体性能。CEP90是案例位置误差的经验90\%分位数。}

\begin{tabular}{lrrrrr}
\toprule
方法 & 均值/m & 中位/m & CEP90/m & RMSE/m & 平均工作量\\
\midrule
静态WLS & 6.532 & 5.340 & 14.959 & 8.504 & 19.3\\
静态稳健 & 6.196 & 3.928 & 14.993 & 8.395 & 20.6\\
恒速稳健 & 4.674 & 2.993 & 9.853 & 6.781 & 27.2\\
旧BIC门控 & 4.347 & 2.142 & 12.203 & 6.790 & 47.8\\
智能体所选 & 4.189 & 2.252 & 9.966 & 6.434 & 71.3\\
\bottomrule\end{tabular}\end{table}

\begin{table}[H]\centering\small
\caption{场景内配对bootstrap比较：对照误差减去所选方法误差。正值表示所选方法较低。}

\begin{tabular}{lrr}
\toprule
对照 & 均值减少/m & 95\%区间/m\\
\midrule
静态WLS & 2.343 & [1.890, 2.776]\\
静态稳健 & 2.007 & [1.578, 2.424]\\
恒速稳健 & 0.485 & [0.064, 0.912]\\
旧BIC门控 & 0.158 & [-0.091, 0.425]\\
\bottomrule\end{tabular}\end{table}

区间描述固定生成机制下的有限样本差异，没有对众多方法比较作多重校正。它不能度量重新运行一次大模型规划器时的方差，也不能据此归因规划器优于其它研究策略。

\subsection{逐场景结果与误差分布}

\begin{table}[H]\centering\small
\caption{四类常规场景的平均当前位置误差（m），每类90例。}

\begin{tabular}{lrrrr}
\toprule
方法 & 静态正常 & 静态异常 & 移动正常 & 移动异常\\
\midrule
静态WLS & 1.031 & 3.492 & 10.349 & 11.259\\
静态稳健 & 1.059 & 2.080 & 10.429 & 11.217\\
恒速稳健 & 2.636 & 6.229 & 2.619 & 7.213\\
旧BIC门控 & 1.305 & 2.080 & 2.984 & 11.019\\
智能体所选 & 1.419 & 2.329 & 3.681 & 9.328\\
\bottomrule\end{tabular}\end{table}

\begin{figure}[H]\centering
\begin{tikzpicture}\begin{axis}[ybar,bar width=9pt,width=0.94\linewidth,height=6.3cm,
ylabel={平均位置误差（m）},ymin=0,grid=major,enlarge x limits=0.15,
xtick={0,1,2,3},xticklabels={静态正常,静态异常,移动正常,移动异常},
legend style={at={(0.5,1.04)},anchor=south,legend columns=3},
x tick label style={font=\small,text width=17mm,align=center},]
\addplot[fill=navy,draw=navy] coordinates {(0,1.030812) (1,3.491518) (2,10.348803) (3,11.258691)};
\addlegendentry{静态WLS}
\addplot[fill=orange,draw=orange] coordinates {(0,2.635838) (1,6.229444) (2,2.619052) (3,7.212711)};
\addlegendentry{恒速稳健}
\addplot[fill=teal,draw=teal] coordinates {(0,1.418681) (1,2.329268) (2,3.681274) (3,9.327780)};
\addlegendentry{智能体所选}
\end{axis}\end{tikzpicture}
\caption{常规场景的配对误差比较。总体改善不能替代逐场景非劣分析。}
\label{fig:main-bars}
\end{figure}

\begin{figure}[H]\centering
\begin{tikzpicture}\begin{axis}[width=0.92\linewidth,height=7cm,xlabel={位置误差（m）},ylabel={经验累计概率},xmin=0,xmax=30,ymin=0,ymax=1,grid=major,legend pos=south east]\addplot[thick,no marks,color=navy] table[x=error,y=cdf] {data/cdf_static_wls.dat};\addlegendentry{静态WLS}
\addplot[thick,no marks,color=orange] table[x=error,y=cdf] {data/cdf_static_robust.dat};\addlegendentry{静态稳健}
\addplot[thick,no marks,color=brick] table[x=error,y=cdf] {data/cdf_dynamic_robust.dat};\addlegendentry{恒速稳健}
\addplot[thick,no marks,color=gray] table[x=error,y=cdf] {data/cdf_legacy_bic.dat};\addlegendentry{旧BIC门控}
\addplot[thick,no marks,color=teal] table[x=error,y=cdf] {data/cdf_llm_selected.dat};\addlegendentry{智能体所选}\end{axis}\end{tikzpicture}
\caption{独立主测试误差的经验分布。横轴显示0--30米；超过30米的样本仍保留在全部统计中。}
\label{fig:cdf}
\end{figure}

\subsection{组件消融与开发前驱}

\begin{table}[H]\centering\small
\caption{冻结方法的控制消融与开发前驱。消融不重新调参。}

\begin{tabular}{lrrrr}
\toprule
方法 & 主测试均值 & 主测试CEP90 & 压力均值 & 与所选均值差\\
\midrule
固定恒速模型 & 4.674 & 9.850 & 8.814 & 0.485\\
固定静态模型 & 6.196 & 14.994 & 12.839 & 2.007\\
开发前驱1 & 4.704 & 12.564 & 12.396 & 0.515\\
\bottomrule\end{tabular}\end{table}

某个消融若优于冻结方法，应作为预先冻结方案的负面证据报告，不能用该测试结果悄悄替换最终方法。未触发的消融表示冻结配置本来没有该组件，避免重复计算伪造独立对照。

\subsection{分布外压力测试}

每类压力场景独立生成40例，改变几何、运动、观测数量、标定或污染条件。下表同时保留均值与尾部误差，以检验常规场景的改善能否延续到这些变化。

\begin{table}[H]\centering\small
\caption{压力场景的均值／CEP90（m），每类40例。}

\begin{tabular}{lrrr}
\toprule
场景 & 静态WLS & 恒速稳健 & 智能体所选\\
\midrule
狭窄几何 & 14.503 / 21.270 & 9.817 / 24.461 & 9.432 / 22.286\\
加速运动 & 10.846 / 15.653 & 4.414 / 7.723 & 5.313 / 11.554\\
稀疏观测 & 10.805 / 17.493 & 5.976 / 14.303 & 11.121 / 18.090\\
标定漂移 & 15.044 / 21.988 & 11.221 / 21.014 & 13.278 / 24.098\\
高比例污染 & 10.862 / 21.049 & 12.641 / 25.044 & 11.734 / 22.187\\
\bottomrule\end{tabular}\end{table}

\begin{figure}[H]\centering
\begin{tikzpicture}\begin{axis}[ybar,bar width=9pt,width=0.94\linewidth,height=6.3cm,
ylabel={平均位置误差（m）},ymin=0,grid=major,enlarge x limits=0.15,
xtick={0,1,2,3,4},xticklabels={狭窄几何,加速运动,稀疏观测,标定漂移,高比例污染},
legend style={at={(0.5,1.04)},anchor=south,legend columns=3},
x tick label style={font=\small,text width=17mm,align=center},]
\addplot[fill=navy,draw=navy] coordinates {(0,14.502767) (1,10.845790) (2,10.804917) (3,15.044343) (4,10.861705)};
\addlegendentry{静态WLS}
\addplot[fill=orange,draw=orange] coordinates {(0,9.817206) (1,4.414207) (2,5.976103) (3,11.221198) (4,12.640590)};
\addlegendentry{恒速稳健}
\addplot[fill=teal,draw=teal] coordinates {(0,9.431836) (1,5.312884) (2,11.121400) (3,13.277567) (4,11.734202)};
\addlegendentry{智能体所选}
\end{axis}\end{tikzpicture}
\caption{冻结估计器在五类压力条件下的误差。压力输入不参与方法选择。}
\label{fig:stress-bars}
\end{figure}

\begin{table}[H]\centering\small
\caption{所选方法实际使用的运动模型次数。标签来自求解器诊断。}

\begin{tabular}{lrrr}
\toprule
场景 & 静态 & 恒速 & 其它\\
\midrule
静态正常 & 78 & 12 & 0\\
静态异常 & 85 & 5 & 0\\
移动正常 & 12 & 78 & 0\\
移动异常 & 49 & 41 & 0\\
狭窄几何 & 11 & 29 & 0\\
加速运动 & 5 & 35 & 0\\
稀疏观测 & 40 & 0 & 0\\
标定漂移 & 21 & 19 & 0\\
高比例污染 & 30 & 10 & 0\\
\bottomrule\end{tabular}\end{table}

\subsection{独立校准、覆盖与拒绝}

使用240个独立混合场景案例校准尺度，再在320个同分布案例及200个压力案例上检查覆盖。名义水平来自预定协议，实测覆盖与区域大小均由已保存的逐例结果计算。

\begin{table}[H]\centering\small
\caption{独立校准后的实际覆盖。压力分布的覆盖没有交换性保证。}

\begin{tabular}{lrrrrrr}
\toprule
名义覆盖 & 分位数$q$ & 评估总体 & 样本 & 实际覆盖 & 均值半径/m & 中位半径/m\\
\midrule
90.0\% & 3.548 & 同分布混合 & 320 & 85.6\% & 8.481 & 7.084\\
90.0\% & 3.548 & 压力分布 & 200 & 74.0\% & 20.147 & 14.300\\
95.0\% & 5.199 & 同分布混合 & 320 & 94.7\% & 12.429 & 10.383\\
95.0\% & 5.199 & 压力分布 & 200 & 85.5\% & 29.527 & 20.958\\
\bottomrule\end{tabular}\end{table}

\begin{table}[H]\centering\small
\caption{90\%名义区域的拒绝曲线：只接受半径不超过阈值的案例。}

\begin{tabular}{lrrrr}
\toprule
半径阈值/m & 接受数 & 接受率 & 接受均值误差/m & 接受后覆盖\\
\midrule
2 & 0 & 0.0\% & -- & --\\
5 & 85 & 26.6\% & 2.034 & 88.2\%\\
10 & 229 & 71.6\% & 3.165 & 83.8\%\\
15 & 291 & 90.9\% & 3.540 & 84.5\%\\
20 & 311 & 97.2\% & 3.702 & 85.2\%\\
30 & 318 & 99.4\% & 3.870 & 85.5\%\\
50 & 320 & 100.0\% & 3.955 & 85.6\%\\
\bottomrule\end{tabular}\end{table}

拒绝阈值越小通常会减少输出数量。表中的接受后覆盖是实际样本统计，不继承总体边际覆盖定理；缺失的接受子集以空值表示。较大的预测圆盘即使覆盖充分，也可能缺乏实际定位价值，因此必须同时报告半径与接受率。

\subsection{独立数学见证}

在标量检验中，真距离为20米，对数噪声标准差为0.35。解析自加权最优值为15.65409米，200000次独立模拟得到15.67145米；修正后对数目标给出20.01864米。该见证用于核对解析公式，不构成新的泛化实验。

镜像构造的两条轨迹当前位置相差6.0米，所有距离差的最大绝对值为0.0e+00米；数据Jacobian秩为4，最小奇异值为0.96509。因此，这个实例直接展示了局部满秩与全局镜像歧义能够共存。

\subsection{哪些结论成立，哪些预期没有泛化}

\paragraph{总体均值有改善，不能据此宣称处处更优。}
冻结的距离域交叉验证方法在360例主测试上的平均误差为4.189米，相比WLS的6.532米降低35.87\%。它比固定恒速稳健基线的均值低10.38\%，但CEP90为9.966米，略高于后者的9.853米。相对旧BIC门控，均值只减少0.158米，配对95\%区间为$[-0.091,0.425]$米，包含零；本轮不能确认其稳定优于BIC。一个总体均值无法同时概括静态、动态与尾部风险。

\paragraph{静止非劣条件没有从验证集推广到测试集。}
验证集静态正常误差为0.902米，低于同批WLS的0.923米；但独立测试中变为1.419米，WLS为1.031米，退化37.63\%。因此，即使把开发判据放宽为10\%容忍度，测试结果仍未满足。更不能把总体35.87\%的均值改善写成完成原题的“标准场景不劣化”或“移动检测精度提高20\%”。原题的后一指标没有得到定义，也没有真实工业对照数据。

\paragraph{门控节省了静态过拟合的代价，也丢失了一部分运动收益。}
固定恒速消融的主测试均值为4.674米，高于所选方案，但其压力测试均值为8.814米，低于所选方案的10.176米。所选方案在稀疏场景的平均误差为11.121米，固定恒速基线为5.976米。这里有可核查的实现机制：三折划分下，8份观测产生不足6条的训练折，运动模型因预设可用性门槛被排除。它不是从数据中“发现目标静止”；这种保守回退需要作为系统失效方式报告。

\paragraph{计算成本提高，稳健性仍有明确边界。}
主测试上所选方法平均消耗71.3个工作单位，旧BIC为47.8、恒速稳健为27.2；本机记录的平均用时约为14.9、7.6和4.9毫秒。工作单位在新旧求解器之间的计数口径有差异，壁钟时间也依赖实现和设备，所以这些数字只说明本次运行的精度与成本取舍。最终配置没有使用截尾或子集初始化，不能把本轮效果归因于这些虽然已实现、但没有进入冻结方案的组件。

\paragraph{不确定性层必须报告实测欠覆盖。}
名义90\%区域在独立同分布混合样本上仅覆盖274/320例，即85.625\%；名义95\%区域覆盖303/320例，即94.6875\%。移动异常子群的相应覆盖为59.76\%和87.80\%，静态两个子群则均为100\%。这展示了总体校准与子群均衡之间的差别。压力分布上的覆盖进一步降至74.0\%和85.5\%，尤其稀疏场景分别只有37.5\%和60.0\%。

共形定理对随机校准样本与新案例联合取概率，不承诺任一已实现校准集下的有限测试比例必定超过名义水平。本次90\%经验覆盖没有达标，这一事实保留；既不能将它掩盖为“90\%覆盖已获验证”，也不能仅凭一个有限样本比例否定带有条件的边际覆盖定理。半径10米的90\%区域拒绝策略在同分布测试中接受71.56\%案例，接受后的覆盖只有83.84\%，再次说明按半径拒绝不会自动保留条件覆盖。

上述结果支持一个有限结论：本研究得到一个对常规混合场景有平均收益的可审计方案，并定位了静止泛化、稀疏门控、尾部风险与子群校准方面的失败。开放问题仍然存在，失败结果本身为下一轮建模提供了更明确的约束。


\subsection{校准随机性的事后数学参考}
\label{sec:coverage-variability-reference}

名义90\%区域在本次独立覆盖检验中的经验覆盖率为
$274/320=85.625\%$，这一不足必须保留。以下计算只解释有限校准样本的随机性，
不是新增实验、事后重新校准或方法修订，也不把未达名义水平改写成成功。

假设位置估计器与正尺度函数均已在校准前冻结，校准分数和测试分数独立同分布，
且其共同累积分布函数$F$连续。令校准样本数$m=240$，
$q=S_{(k)}$，$k=\lceil(m+1)(1-\alpha)\rceil$。
概率积分变换与次序统计量给出
\[
 U=F(q)\sim\operatorname{Beta}(k,m+1-k),\qquad
 K\mid U\sim\operatorname{Binomial}(320,U),
\]
其中$K$为320个新案例中被区域覆盖的数量。对随机$U$积分后，
得到
\[
 K\sim\operatorname{BetaBinomial}(320,k,m+1-k).
\]
这是\emph{校准集和测试集共同重新抽样}的无条件分布；
固定本次校准集后的覆盖概率$F(q)$仍未知，因此下表不是该条件覆盖率的置信区间。

\begin{table}[H]
\centering\small
\caption{理想连续IID评分模型下的事后覆盖计数参考。预测区间取离散分布的2.5\%和97.5\%分位数。}
\label{tab:coverage-variability-reference}
\begin{tabular}{@{}crrrcr@{}}
\toprule
名义覆盖 & $k$ & 实测$K$ & $\mathbb{E}[K]$ & 95\%预测区间 & $\Pr(K\le K_{\rm obs})$\\
\midrule
90\% & 217 & 274 & 288.13 & $[271,303]$ & 0.0563\\
95\% & 229 & 303 & 304.07 & $[291,314]$ & 0.4258\\
\bottomrule
\end{tabular}
\end{table}

对应的覆盖比例预测区间分别为$[84.6875\%,94.6875\%]$与
$[90.9375\%,98.125\%]$。数值使用\texttt{scipy.stats.betabinom}计算，
参数、软件版本及完整小数保存在
\nolinkurl{research/coverage_variability_reference.json}。
这些参考说明本次较低覆盖与理想联合抽样模型并不必然矛盾，但不能验证实际
评分分布的连续IID假设、保证给定校准集的覆盖，或弥补已观察到的85.625\%覆盖率。
它们也不提供分场景、拒绝子集或分布偏移下的保证；研究未据此调整$q$、
尺度函数、阈值或定位方法。
